% !TEX root = ../../../main.tex
% Block card: catalogue.multiscale_entropy -- "Multiscale entropy" (Part I, the blocks chapter).
\cardfield{What it does}
Draws a ``complexity curve'' for each long window: how unpredictable the signal is when you
look at it in finer or coarser detail.

\begin{plainterms}[Multiscale entropy]
Read a signal one value at a time and ask how surprised you are by the next one: that is its
entropy. Now blur the signal by averaging every 2 values, then every 10, then every few
hundred, and ask again each time. Pure hiss stops surprising you once blurred, because the
averages all look alike. A signal with structure at many time-scales keeps surprising you
however much you blur it. The curve of surprise against blur is the fingerprint
\citep{costa2005multiscale}.
\end{plainterms}

\cardfield{When to use it}
To compare recordings, species or conditions by their complexity, not to find events: a
curve needs hours of signal. The \emph{complexity\_fingerprint} template is the usual chain:
a 2-hour detrend, 72-hour windows, then this block.

\cardfield{What goes in, what comes out}
In goes a window set on one channel, windows hours long. Out comes the same window set with
each window's curve (one number per scale) and a few summaries of it added to its measures.

\cardfield{Settings that matter}
\emph{scale\_max\_s} is the coarsest blur. A window only reaches about its length divided by
1{,}000, so a 24-hour window reaches about 86\,s; scales beyond that are left blank.
\emph{r\_frac} sets how close two stretches must be to count as alike. \emph{r\_scale} chooses
how the window's spread is measured; the default, \emph{mad}, is less thrown by outliers.

\cardfield{Watch out for}
Remove slow drift first (\blockref{preprocessing.detrend}): drift makes every scale look
regular, and the block flags windows where that happened. Coloured noise has a curve of its
own, so judge a curve against amplitude-adjusted surrogate copies of the same recording
(\blockref{preprocessing.surrogate}, method \emph{aaft}); the block does not do that
comparison for you.
